\section{Random restriction to an accurate face}\label{sec:face-restriction}
For a random $m_f$-face and its degree-$u$ Fourier truncation,
\begin{equation}\label{eq:face-restriction}
 \E_{\rm face}\|P-P_{\rm face,\le u}\|_2^2
 \le\|P\|_2^2
       \left(\frac{e(r/k)m_f}{u+1}\right)^{u+1}.
\end{equation}
To see this, expand in Walsh characters. A support of size at most $r$
retains more than $u$ active coordinates with at most the displayed
hypergeometric upper-tail probability. Random fixed signs cancel all
cross terms. Averaging and summing the squared coefficients proves the
inequality.

If the global head error is at most $3/8+o(1)$, a positive fraction of
faces have error at most $7/16$ and hence correlation at least $1/8$.
Each sign class on such a face has mass at least $1/16$.
The global head leakage and Markov discard only a vanishing fraction
for the bad-margin eligibility condition. The population versions of
\eqref{eq:weighted-empirical-cost}--\eqref{eq:weighted-empirical-huge},
together with~\eqref{eq:face-restriction}, with fixed
Markov constants, select one accurate eligible face with controlled
wrong-side moment, huge-set mass, and restriction leakage.
Normalize its truncation with the same $s,t$:
$P_0=(P_{\rm face,\le u}-t)/s$. The contact projection bound gives
\[
 \|F-P_0\|_p\le C\sqrt{\zeta_0}
                +2\|P-P_{\rm face,\le u}\|_2/a.
\]
Take $u=a_1H$ with $a_1$ sufficiently small relative to $b_0$.
Then the huge-set mass $C\rho^{b_0}$ satisfies the absorption condition
in Theorem~\ref{thm:endpoint-contact-calibration}. Its conclusion forces
$u\ge c m_f^\theta$.
